% !PRINTED 269
\weberpage{269}
%<<CONT
Zeichenwechsel, bei negativer Discriminante, wo auch $a$ negativ ist, \gesperrt{ein} Zeichenwechsel stattfindet.

Unsere Gleichung hat also, so lange $a$ in dem Intervall (9) liegt, zwei imaginäre Wurzeln, wenn es ausserhalb dieses Intervalles liegt, vier imaginäre Wurzeln, niemals vier reelle Wurzeln.

Ein rein numerisches Beispiel bietet die der complexen Multiplication der elliptischen Functionen entnommene Gleichung
\[
x^5 - x^3 - 2x^2 - 2x - 1 = 0.
\]
Für die Determinante der Bezoutiante erhalten wir aus \S.\,73
\[
\begin{vmatrix}
\dfrac{-4}{5}, & \dfrac{-3}{5}, & \dfrac{4}{5}, & 5 \\[6pt]
\dfrac{-3}{5}, & \dfrac{4}{5}, & \dfrac{33}{5}, & 8 \\[6pt]
\dfrac{4}{5}, & \dfrac{33}{5}, & \dfrac{46}{5}, & 6 \\[6pt]
5, & 8, & 6, & 2
\end{vmatrix}
\]
die Kette der Haupt-Unterdeterminanten
\[
\frac{47^2}{5}, \; \frac{94}{5}, \; -1, \; \frac{-4}{5}, \; 1,
\]
so dass die Discriminante unserer Gleichung $47^2$ ist. Die Gleichung hat also zwei Paar imaginäre und eine reelle Wurzel. Das zweite Glied dieser Kette $94 : 5$ braucht man, wenn man die Discriminante kennt, nicht zu berechnen, ja es genügt schon die Kenntniss des negativen Vorzeichens im vorletzten Gliede, um die vorstehende Entscheidung über die Realität der Wurzeln zu treffen.
